\begin{textsample}{POS Dim 2 – am – Score 23.00 – am/am\_ef\_19.txt}  \label{ex:f2_pos_157}
6.
A formação do professor A Educação Especial na perspectiva da Educação Inclusiva é uma temática bem discutida na sociedade contemporânea.
Reforça e objetiva incluir todos no âmbito escolar e na sociedade e, com isso, fazer cumprir o que garante a Constituição Federal Brasileira, que afirma que todos têm direito à educação.
Torna -se indispensável que a formação do professor da Educação Especial na perspectiva da Educação Inclusiva seja uma abordagem objetiva e reflita sobre uma melhor compreensão dessa modalidade de ensino.
Uma forma de compreender melhor esse modelo de ensino é conhecer um pouco da história e das \textbf{legislações} que tratam sobre essa modalidade.
Identificar também como é realizada a formação do professor que atua de forma direta nesse viés, buscando encontrar em sua formação teórico-prática elementos que proporcionem a inclusão do público-alvo da Educação Especial nas escolas de ensino regular e, consequentemente, na sociedade. 6.1.
Historicamente, as pessoas com deficiência sofreram com várias formas de preconceito que se estenderam durante muitos séculos e, apesar de hoje existirem leis que \textbf{amparam} seus direitos enquanto cidadãos, ainda se faz necessário que outras medidas sejam tomadas para impedir ou até amenizar algumas formas de exclusão que ainda são imperativas em nossa sociedade.
Com os avanços dos Direitos Humanos, registraram -se consideráveis progressos na conquista da igualdade e do exercício de direitos.
É o que se sente e o que se observa atualmente, tendo como grande enfoque a busca da inclusão destas pessoas historicamente marcadas pela segregação, pelo preconceito e pela rejeição.
Nesse processo, é de suma importância que haja um trabalho conjunto com toda a comunidade escolar, fazendo com que estejam todos cientes e envolvidos nesta causa, tornando o ambiente acessível para essas pessoas.
Um ambiente educacional acessível e favorável a esses educandos configura um indicativo de que a inclusão pode de fato acontecer, e não somente a integração.
De acordo com Sassaki ( 2006 ), a integração propõe a inserção parcial do sujeito, enquanto que a inclusão propõe a inserção total.
Para que isso aconteça, a escola precisa romper com a perspectiva homogeneizadora e adotar estratégias para assegurar os direitos de aprendizagem de todos.
Segundo esse mesmo autor, as escolas precisam romper com diversas barreiras em vários aspectos de suas estruturas tradicionais para favorecer a acessibilidade no atendimento desses alunos.
De acordo com Sassaki ( 2009 ), são seis as principais barreiras a serem consideradas com total atenção : > [... ] arquitetônica ( sem barreiras físicas ), comunicacional ( sem barreiras na comunicação entre pessoas ), metodológica ( sem barreiras nos métodos e técnicas de lazer, trabalho, educação etc. ), instrumental ( sem barreiras nos instrumentos, ferramentas, utensílios etc. ), programática ( sem barreiras embutidas em \textbf{políticas}, \textbf{legislações}, normas etc. ) e atitudinal ( sem preconceitos, estereótipos, estigmas e discriminações nos comportamentos da sociedade para a pessoa que tem deficiência ). ( SASSAKI, 2009, p. 1-2 ) Nesse cenário, é preciso o envolvimento de vários \textbf{profissionais} de outras áreas de conhecimento, que atuarão em conjunto com os \textbf{profissionais} da educação, como os \textbf{profissionais} da saúde, fonoaudiólogos, psicólogos, assistentes sociais e, principalmente, a participação da família, por esta estabelecer um contato direto com a pessoa com deficiência.
Sendo assim, precisa estar bem informada a respeito da condição de seu assegurado e apta aos cuidados de seu familiar com deficiência, consolidando, assim, um trabalho interdisciplinar de \textbf{qualidade} e integral para com os diferentes.
O professor que atua de forma direta realiza seu trabalho pedagógico no ambiente escolar, busca promover as potencialidades dos sujeitos envolvidos no processo de ensino-aprendizagem e torna -se um \textbf{profissional} relevante e insubstituível nessa modalidade de educação.
A esse respeito, corrobora Cury ( 2003 ) : > Os educadores, apesar de suas dificuldades, são insubstituíveis porque a gentileza, a solidariedade, a tolerância, a inclusão, os sentimentos altruístas, enfim, todas as áreas da sensibilidade não podem ser ensinadas por máquinas, e sim por seres humanos.
A educação, como um fenômeno de responsabilidade social, possui um papel importante para o desenvolvimento e crescimento de uma nação que não pode ser ignorado por nenhum cidadão, sobretudo pelo poder público, pois é a escola que prepara os indivíduos para sua atuação e para exercer sua plena cidadania, gozando de todos os bens e serviços que sejam oferecidos por essa sociedade da qual faz parte.
Para que os resultados educacionais das pessoas com deficiência sejam alcançados com sucesso e, ainda, recebam uma formação adequada, faz -se necessário que os educadores envolvidos nesse processo, como gestores, pedagogos e professores, recebam uma formação de \textbf{qualidade} que possa contribuir de forma positiva para esses resultados e que estes \textbf{atendam} aos objetivos da sociedade.
A \textbf{qualidade} da formação desses \textbf{profissionais}, tanto inicial quanto continuada, é de total importância para refletir na \textbf{qualidade} do ensino de seus discentes.
A UNESCO, em seu relatório de 1998, relacionou a \textbf{qualidade} do ensino com a formação qualificada desses \textbf{profissionais} ( UNESCO, 1998 apud Roseli Rocha de Carvalho, Educação Especial : do querer ao fazer, p. 28 ) : > A \textbf{qualidade} dos serviços educacionais para pessoas com deficiência depende da \textbf{qualidade} da formação.
Esta deverá ser parte integrante dos planos nacionais, onde se contemplam os \textbf{requisitos} dessa formação.
Houve um momento na história do nosso país em que a formação para os professores era feita de uma forma \textbf{técnica}, apenas pelo magistério, em que o aluno do ensino médio era direcionado para escolas \textbf{técnicas} para receber uma formação direcionada a uma atuação como professor nas séries iniciais.
Os estudos para essa modalidade eram pautados nas metodologias de ensino, domínio dos conteúdos e dinâmica da turma, ou seja, nas práticas realizadas nas salas de aula.
Com a expansão do ensino no território nacional, surgiram vários problemas relacionados à \textbf{qualidade} do ensino e da aprendizagem.
Na década de noventa, exigia -se do \textbf{profissional} da educação uma nova postura, fundamentada na teoria construtivista interacionista concebida por Jean Piaget ( 1896-1980 ).
Para esse teórico, a criança, como um ser biológico, é ativa e age espontaneamente no meio, ou seja, é pelo contato com o mundo que seus conhecimentos são construídos.
Destaca Furtado ( 2009, p. 139 ) : > O ser humano, dotado de estrutura biológica, herdada por uma forma de funcionamento intelectual, ou seja, uma maneira de interagir com o ambiente que o leva à construção de um conjunto permitirá a organização desses significados em estruturas cognitivas.
É importante que o professor possua conhecimentos dessa natureza, que contempla a psicologia do desenvolvimento humano, disciplina essa que compõe a grade curricular do \textbf{curso} de formação inicial de professores.
Nesse sentido, para poder conhecer melhor seus educandos, através desse conhecimento e de um olhar cuidadoso, fica mais fácil para o professor identificar algumas deficiências ou dificuldades de aprendizagem, como dislexia, dislalia, deficiência intelectual, transtornos do espectro autista e outros.
Foi com o intuito de reparar e melhorar o atendimento educacional nas escolas brasileiras que houve uma exigência por parte do poder público \textbf{vigente}, a qual resultou na aprovação, em 2009, de um projeto que torna obrigatório que todos os professores do ensino básico tenham diploma universitário e licenciatura.
A Lei de \textbf{Diretrizes} e Bases da Educação ( LDB ) determina que todos os professores devam possuir \textbf{curso} superior.
A partir de então, o \textbf{curso} de Pedagogia passou a formar esses \textbf{profissionais} para atuar na Educação Infantil e nos Anos Iniciais.
São nessas etapas de ensino que inicialmente foram direcionados os estágios supervisionados, pois, neste primeiro momento, não constava nessa formação estudo para a Educação Especial.
Todavia, com a assinatura da Declaração de Salamanca, em 1994, o Brasil passou a rever o atendimento de crianças com deficiências, o que levantou uma questão a respeito da formação de professores, pois, no entendimento de alguns educadores, não era obrigatório formar professores especializados nessa modalidade de ensino, pois a inclusão passaria a ser tarefa de todas.
Diante dos desafios da Educação Especial, surge a importância do aperfeiçoamento dos professores para o exercício de serviços educacionais especializados.
A \textbf{Portaria} nº 1.793, de dezembro de 1994, recomendou a implementação, nos \textbf{currículos} de formação docente, da disciplina “ Aspectos ético-político educacionais da \textbf{normatização} e integração da pessoa com necessidades especiais ”.
Indicou a manutenção e expansão dos estudos como \textbf{cursos} e \textbf{especialização} na área desse conhecimento, com o objetivo de melhorar a \textbf{qualidade} do atendimento para esta modalidade.
Também consta na Lei de \textbf{Diretrizes} e Bases da Educação Nacional ( LDBN 9.394/96 ) que os professores tenham \textbf{especialização} adequada em nível médio ou superior para o atendimento especializado, bem como professores do ensino regular capacitados para a integração desses educandos nas classes comuns.
Consta no Plano Nacional de Educação ( PNE ), com vigência de 2014 a 2024, que seus eixos temáticos contemplem em suas \textbf{diretrizes}, metas e estratégias de ações um melhor atendimento nas escolas de todo o país : I ) Papel do Estado na \textbf{garantia} do direito à educação de \textbf{qualidade} : organização e \textbf{regulamentação} da educação nacional ; II ) Qualidade da educação, gestão democrática e avaliação ; III ) Democratização do acesso, permanência e sucesso escolar ; IV ) Formação e valorização dos/das \textbf{profissionais} da educação ; V ) Financiamento da educação e controle social ; VI ) Justiça social, educação e trabalho ; inclusão, diversidade e igualdade ( BRASIL, 2015 ).
Todos esses elementos legais vêm corroborando no sentido de tornar o sistema educacional brasileiro mais significativo e eficaz e, com isso, contribuindo com a \textbf{qualidade} da formação dos professores, conforme asseveram as metas 15 e 16 do referido documento : > [... ] assegurado que todos os professores e as professoras da educação básica possuam formação específica de nível superior, obtida em \textbf{curso} de licenciatura na área de conhecimento em que atuam. > > [... ] formar em nível de pós-graduação, 50 \% ( cinquenta por cento ) dos professores da educação básica, até o último ano de vigência deste PNE, e garantir a todos ( as ) os \textbf{profissionais} da educação básica formação continuada em sua área de atuação, considerando as necessidades, demanda e contextualizações de ensino. ( BRASIL, 2015, p. 79-80 ).
Entretanto, a educação das pessoas com deficiência, que no início estabelece o desenvolvimento de novas ações pedagógicas na intenção de explorar e potencializar habilidades e competências, foi mal interpretada.
Por falta de adequação ou incapacidade de seus educandos, passou a segregar mais do que incluir na interação social ( PACHECO ; ALVES, 2017 ).
Na tradição, o \textbf{curso} de Pedagogia não era \textbf{destinado} a formar professores, mas sim a formar educadores, planejadores, gestores e pesquisadores como bacharelado.
Depois, foi agregado mais um ano de didática e práticas de ensino para formar professor, mas não alfabetizar, e sim para dar aula nas escolas normais.
O processo de formação de professores ainda carece de maiores ajustes para se tornar um \textbf{curso} que forme de fato \textbf{profissionais} que atuem nas salas de aula e saibam desenvolver um trabalho junto a um público diversificado.
E por mais que ocorra a inserção de métodos singulares na construção da formação do educador, como já abordado, a percepção do docente pode não condizer com sua visão de Educação Especial trazida por essas práticas.
Por isso, é premente ressaltar o papel que cabe à formação continuada, que ocupou um espaço considerável, onde \textbf{profissionais} da educação têm a oportunidade de expandir e agregar novos conhecimentos, proporcionando ao docente em atividade uma excelente oportunidade para as trocas de experiência com outros \textbf{profissionais} da área, tornando -os mais qualificados para assumirem uma melhor postura diante de seus alunos.
Dessa maneira, infere Baumel ( 2003 ) : > A formação de professores deve ser concebida como um contínuo, associado à compreensão do desenvolvimento \textbf{profissional} ; em outras palavras, formar e articular “ uma variedade de formatos de aprendizagem ”.
O comprometimento, aqui, é de interligar a formação inicial com a continuada, que não abarca o termo e o processo de \textbf{capacitação}.
O processo de formação inicial e continuada é um projeto diferenciado, em fases, ao longo de uma finalidade e um estado de desenvolvimento \textbf{profissional} ( BAUMEL, 2003, p. 30 ).
Tratando -se de Educação Especial na perspectiva da educação inclusiva, os desafios são constantes e requerem desse \textbf{profissional} controle e domínio das diversas situações do cotidiano escolar voltadas a essa modalidade.
O professor pode encontrar essa formação no próprio espaço da escola e nos núcleos de apoio \textbf{ofertados} pelas \textbf{redes} de ensino nas quais estiver inserido, por meio de palestras, \textbf{oficinas}, leituras e \textbf{cursos}.
O mesmo deve ter a vontade de ir em busca de novos conhecimentos, no sentido de avançar com um \textbf{currículo} diferenciado em relação à \textbf{qualidade} de suas práticas educativas com seus educandos.
A esse respeito, aduz Prieto ( 2006 ) : > A formação continuada do professor deve ser um compromisso dos sistemas de ensino comprometidos com a \textbf{qualidade} do ensino que, nessa perspectiva, devem assegurar que sejam aptos a elaborar e a implantar novas propostas e práticas de ensino para responder às características de seus alunos, incluindo aquelas evidenciadas pelos alunos com necessidades educacionais especiais ( PRIETO, 2006, p. 57 ).
É por meio dessa formação continuada que o professor pode ampliar seu olhar para contemplar as potencialidades de seus alunos e não enxergar somente suas limitações.
Diante dos desafios que se referem à Educação Especial e Inclusiva, é indispensável que o professor mantenha uma constante formação teórico-prática, a fim de \textbf{atender} melhor nesse ambiente tão diversificado que é o chão das escolas.

% matched lemmas: amparar, atender, capacitação, currículo, curso, destinar, diretriz, especialização, garantia, legislação, normatização, ofertar, oficina, política, portaria, profissional, qualidade, rede, regulamentação, requisito, técnico, vigente
\end{textsample}
